\documentclass[11pt]{article}
\usepackage{kristiania-workbook}

\begin{document}
\begin{center}
  {\LARGE\bfseries Put a web service under load}\\[3pt]
  {\large Session 3 -- Web services, load \& DDoS basics}\\[5pt]
  {\small Exercise 3 \textperiodcentered\ Session 3 lab \textperiodcentered\ Estimated time: 75--100 min}
\end{center}
\vspace{2pt}\hrule\vspace{1.1em}

Write your answers in the answer document \cmd{SKY2100\_Exercise3\_Answers.docx}. Today you stand up a real web server, attack it with
\emph{controlled} load, and read what the metrics and logs show. You will watch a service degrade
gracefully and then hit the wall -- the same shape whether the load is real customers or a
denial-of-service attack. This is your first hands-on encounter with the \textbf{Availability} leg of
the CIA triad.

\caution{Run load \emph{against your own VM only} (\cmd{localhost}). Pointing a load generator at a
machine or site you do not own is an attack, and it is illegal.}

\wbsection{1\quad Stand up a web server}
\task{Install Nginx and the benchmark tool:} \cmd{sudo apt update \&\& sudo apt install -y nginx apache2-utils}. Open \cmd{http://localhost/} in the VM browser and confirm the welcome page loads.

\wbsection{2\quad Generate and watch load}
\task{Watch the system:} in terminal~1 run \cmd{htop} (or \cmd{btop}); keep it visible.

\task{Generate load} (terminal~2): \cmd{ab -n 2000 -c 100 http://localhost/}. Record in the answer
document: requests per second, time per request (mean), failed requests.

\task{Increase the pressure:} repeat with \cmd{ab -n 5000 -c 300 http://localhost/}. Note at roughly how
many concurrent requests response time climbs sharply.

\task{Read the logs:} \cmd{sudo tail -n 20 /var/log/nginx/access.log}; correlate the burst of log
lines with the spike you saw in \cmd{htop}.

\aha{Watch what happened as you raised the concurrency. Requests per second climbed for a while, then
flattened, and past a certain point response times shot up and requests began to fail. That is your
VM meeting its capacity ceiling, and the shape of the curve is the same whether the load is real
customers on a launch day or an attacker sending a flood. This is why availability is a security
property, not just an operations concern: a denial-of-service attack does nothing clever, it simply
pushes you past the point you just found. Note the concurrency where your own server broke, because
every defence later in the course, from rate limiting to autoscaling, is about moving that point or
soaking up load before it arrives.}

\task{Interpret your run in 2--3 sentences: what happened to CPU/RAM and requests-per-second as
concurrency rose?}

\wbsection{Optional stretch -- a first control}
Add a rate limit (\cmd{limit\_req\_zone} / \cmd{limit\_req}), reload Nginx, re-run the high-concurrency
test, and compare the failed-request count before/after.

\wbsection{3\quad Cloud transfer: Microsoft Learn}
\task{Complete ``Describe Azure compute services'' and ``Describe Azure networking services''}
(two modules in Part~2; the old combined module was split in 2026, and the compute module still lives
at the old combined address).
\par\noindent\mbox{\href{https://learn.microsoft.com/training/modules/describe-azure-compute-networking-services/}{learn.microsoft.com/training/modules/describe-azure-compute-networking-services}}
\par\noindent\mbox{\href{https://learn.microsoft.com/training/modules/describe-azure-networking-services/}{learn.microsoft.com/training/modules/describe-azure-networking-services}}

\task{Map today's ideas to Azure services} (one line each in the answer document): Load Balancer (L4),
Application Gateway (L7 + WAF), Traffic Manager / Front Door (global), Azure DDoS Protection.

\wbsection{4\quad Azure reflection}
\task{Your local \cmd{ab} test hit one VM's ceiling. Which Azure service would let the \emph{same}
application survive the same load, and what does it trade away (cost, complexity)?}

\task{A \textbf{composite SLA}: if a solution chains a web app (99.95\%) behind a gateway (99.95\%),
is the combined availability higher or lower than 99.95\%? Explain in one line.}

\wbsection{5\quad Knowledge check}
\begin{enumerate}
\item A DDoS attack primarily targets which CIA property?\\
      \cbox Confidentiality \quad \cbox Integrity \quad \cbox Availability
\hint{It floods the service offline — which of the three properties is that?}
\item Name the three flavours of DDoS and give one example of each.
\hint{By raw volume, by protocol tricks, and at the application layer.}
\item \cmd{ab -n 1000 -c 50} means\dots\\
      \cbox 1000 concurrent requests \quad \cbox 1000 requests total, 50 at a time \quad
      \cbox 50 requests total \quad \cbox 1000 servers
\hint{A thousand requests in total, fifty running at a time.}
\item Why is an \textbf{application-layer} DDoS the hardest to tell apart from legitimate traffic?
\hint{Its requests are valid and, one by one, indistinguishable from real users.}
\item An Azure \textbf{Application Gateway} operates at which layer, and which security control does it
      host?
\hint{Layer 7 — and it is where the WAF lives.}
\item What is a \textbf{composite SLA}, and why can adding services \emph{lower} overall availability?
\hint{Chained services multiply their availabilities, so the combined figure drops.}
\item List \emph{two} controls on your own server and \emph{two} that work only at cloud scale,
      against DDoS.
\hint{Local: rate limits, caching, timeouts. Cloud: autoscaling, CDN, DDoS protection.}
\item Define \textbf{capacity} and \textbf{availability}, and state the relationship between them.
\hint{One is work per second the system can do; one is being reachable when needed.}
\item According to Comer (Ch.~1), what early problem led to clusters of servers behind a load
      balancer?
\hint{Early sites outgrew a single machine, so many ran behind a balancer.}
\item Distinguish \textbf{scale-up} from \textbf{scale-out}, and say which the cloud favours.
\hint{A bigger single box versus more boxes — the cloud prefers more.}
\item Multi-tenancy + autoscaling produces a ``moving attack surface''. Explain in one sentence why
      that complicates security.
\hint{Instances come and go and tenants share hardware — there is no fixed list to guard.}
\item (Reflection -- from your data) At what concurrency did response time climb sharply, and what does
      that tell you about \emph{your} VM's capacity?
\end{enumerate}

\signoff

\clearpage
\solheader
{\LARGE\bfseries\color{kred} Solutions}\\[2pt]
{\small Work through the lab first, then check here. Model answers are indicative — equivalent reasoning is fine.}\par\vspace{8pt}
\noindent\textbf{Note:} Model answers are indicative -- accept equivalent reasoning. Multiple-choice
answers are in \textbf{bold}.

\wbsection{Knowledge check}
\begin{enumerate}
\item \textbf{Availability.}
\item \textbf{Volumetric} (UDP flood / amplification), \textbf{Protocol} (SYN flood),
      \textbf{Application-layer} (HTTP flood of an expensive endpoint).
\item \textbf{1000 requests total, 50 at a time.}
\item It uses small, valid-looking requests, individually indistinguishable from real users; only
      aggregate behaviour and logs reveal it.
\item \textbf{Layer 7} (application layer); it hosts the \textbf{WAF}.
\item A \textbf{composite SLA} is the combined SLA of chained services; multiplying availabilities
      (e.g.\ $0.9995\times0.9995$) yields a \emph{lower} number than any component -- so more
      components means lower combined availability.
\item Local: rate limiting, caching, timeouts and connection limits. Cloud-scale: autoscaling,
      CDN/anycast, cloud DDoS protection, WAF.
\item \textbf{Capacity} = work per unit time a system can do; \textbf{Availability} = reachable/
      responsive when needed; availability fails once demand exceeds capacity.
\item Popular early web sites exceeded a single machine's capacity, so operators ran a \emph{cluster}
      of identical servers behind a load balancer. (Comer Ch.~1.5, p.~8, verified.)
\item \textbf{Scale-up} = a bigger single machine; \textbf{scale-out} = more machines. The cloud
      favours \textbf{scale-out}.
\item Instances appear/disappear and tenants share hardware, so you cannot secure a fixed list of
      servers -- only a pattern/policy; the surface changes continuously.
\item Reflection from each group's own measurement -- no fixed value.
\end{enumerate}

\wbsection{References}
Comer Ch.~1 (clusters \& load balancing \textbf{p.~8}, economic motivation \textbf{p.~10}) \& Ch.~2
(multi-tenancy \textbf{p.~15}, elastic computing \textbf{p.~16}, IaaS \textbf{p.~19}, SaaS
\textbf{p.~20}) -- all verified; J\o sang Ch.~2 \& 6 (syllabus, page numbers not verified). MS Learn:
\emph{Describe Azure compute services} and \emph{Describe Azure networking services} (verified, June
2026 -- the previously combined module was split into two).

\end{document}
